% Part: lambda-calculus
% Chapter: lambda-definability
% Section: partial-recursive-functions

\documentclass[../../../include/open-logic-section]{subfiles}

\begin{document}

\olfileid[hi]{lam}{ldf}{par}
\olsection{आंशिक पुनरावर्ती फलन \usetoken{S}{lambda definable} हैं}

आंशिक पुनरावर्ती फलन वे हैं जो मूल फलनों से संयोजन, आदिम पुनरावर्तन और
अपरिबद्ध न्यूनीकरण द्वारा प्राप्त होते हैं। वे सामान्य पुनरावर्ती फलनों से
इस बात में भिन्न हैं कि अपरिबद्ध खोज में उपयोग होने वाले फलनों का नियमित
होना आवश्यक नहीं। नियमितता की माँग न करने का अर्थ है कि न्यूनीकरण द्वारा
परिभाषित फलन कभी-कभी अपरिभाषित हो सकते हैं।

पहली दृष्टि में लग सकता है कि जिन विधियों से दिखाया गया कि सभी (संपूर्ण)
सामान्य पुनरावर्ती फलन !!{lambda definable} हैं, उन्हीं से सभी आंशिक
पुनरावर्ती फलनों को !!{lambda definable} सिद्ध किया जा सकता है। उदाहरणार्थ,
यदि $f$ और $g$ क्रमशः पदों $F$ और~$G$ द्वारा !!{lambda defined} हैं, तो
$f$ का $g$ के साथ संयोजन $\lambd[x][F (G x)]$ द्वारा !!{lambda defined}
है। किंतु फलन आंशिक हों तो यह समस्याजनक है। जब $g(x)$ अपरिभाषित हो, तब $G
x$ का कोई प्रसामान्य रूप नहीं होता। अधिकांश स्थितियों में इसका अर्थ है कि
$F (G x)$ का भी कोई प्रसामान्य रूप नहीं, और हम यही चाहते हैं। किंतु उस
स्थिति पर विचार करें जब $F$, ~$\lambd[x][\lambd[y][y]]$ हो; तब $F (G x)$
का प्रसामान्य रूप ($\lambd[y][y]$) होता है।

यह समस्या दुस्तर नहीं है, और सभी आंशिक पुनरावर्ती फलनों को इस प्रकार
!!{lambda define} करने की विधियाँ हैं कि अपरिभाषित मानों को प्रसामान्य रूप
रहित पदों से निरूपित किया जाए। फिर भी, ये विधियाँ सामान्य पुनरावर्ती फलनों
के लिए अपनाए हमारे दृष्टिकोण से कुछ अधिक जटिल और कम सहज हैं। यहाँ हम प्रमेय
को बिना प्रमाण के दर्ज करते हैं:

\begin{thm}
  सभी आंशिक पुनरावर्ती फलन !!{lambda definable} हैं।
\end{thm}

\end{document}

